\src{190}
which range over the category of coherent algebraic sheaves on $X$ and take
values, respectively, in the category of coherent algebraic sheaves on $X$ and
in the category of modules over $H^{0}(X,\Osh_X)$, by the formulas
\begin{equation}\tag{8.1}
\resizebox{\textwidth}{!}{$
\left\{
\begin{aligned}
\underline{T}_{r}^{s*}(A_1,\ldots,A_r;B_1,\ldots,B_s)
  &={}H^{*}\!\left(
    \Homsh_{\Osh_X}\!\left(
      L(A_1)\otimes\cdots\otimes L(A_r),
      L(B_1)\otimes\cdots\otimes L(B_s)
    \right)\right),\\[3pt]
T_{r}^{s*}(A_1,\ldots,A_r;B_1,\ldots,B_s)
  &={}\mathbf R^{*}\Gamma\!\left(
    \Homsh_{\Osh_X}\!\left(
      L(A_1)\otimes\cdots\otimes L(A_r),
      L(B_1)\otimes\cdots\otimes L(B_s)
    \right)\right).
\end{aligned}
\right.
$}
\end{equation}
Here $L(F)$ denotes a finite locally free resolution of the coherent algebraic
sheaf $F$, and $\mathbf R^{*}\Gamma(K)$ denotes the hypercohomology of $X$ with
respect to the complex of sheaves $K$. If $r$ or $s$ is zero, the tensor product
of the $L(A_i)$, respectively of the $L(B_j)$, is replaced by $\Osh_X$. In
particular, $\underline{T}_{0}^{0}$ and $T_{0}^{0}$ are graded functors with no
arguments: $\underline{T}_{0}^{0}$ is concentrated in degree 0, where it is the
sheaf $\Osh_X$, while $T_{0}^{0}$ is identified with $H^{*}(X,\Osh_X)$. The
independence of the right-hand sides of (8.1) from the chosen resolutions is in
any event evident for the first line (the question then being local), and for
the second follows from the preceding general remarks, taking into account the
hypercohomology spectral sequence of the complex of sheaves
\[
  K=\Homsh_{\Osh_X}\bigl(L(A_1)\otimes\cdots,
      L(B_1)\otimes\cdots\bigr),
\]
which abuts to the hypercohomology of $X$ with respect to $K$. Its initial term
is $H^{p}(X,H^{q}(K))$, i.e.
\begin{equation}\tag{8.2}
  E_{2}^{p,q}
  =H^{p}\!\left(X,
    (\underline{T}_{r}^{s})^{(q)}
      (A_1,\ldots,A_r;B_1,\ldots,B_s)\right).
\end{equation}
It follows that this spectral sequence itself is independent of the chosen
resolutions, and its abutment is the first member of (8.1). The connecting
operators in the various arguments $A_i,B_j$ are readily defined, since every
exact sequence
$0\longrightarrow F'\longrightarrow F\longrightarrow F''\longrightarrow0$
can be resolved by an exact sequence of finite locally free complexes.

On the system of functors $\underline{T}_{r}^{s*}$, respectively $T_{r}^{s*}$,
one defines operations analogous to those of tensor calculus, whose definitions
follow immediately from (8.1). Thus there is a composition (generalizing the
one considered in Section 2)
\begin{equation}\tag{8.3}
\resizebox{\textwidth}{!}{$
\begin{aligned}
&T_{r}^{s*}(A_1,\ldots,A_r;B_1,\ldots,B_s)
\times T_{r'}^{s'*}(A'_1,\ldots,A'_{r'};B'_1,\ldots,B'_{s'})\\
&\qquad\longrightarrow
T_{r+r'}^{s+s'*}
(A_1,\ldots,A_r,A'_1,\ldots,A'_{r'};
 B_1,\ldots,B_s,B'_1,\ldots,B'_{s'}),
\end{aligned}
$}
\end{equation}
satisfying the evident associativity and compatibility properties with respect
to functorial homomorphisms, connecting homomorphisms, and spectral sequences.
There are likewise symmetry operations, whose explicit description is left to
the reader. In addition, there is a contraction operation whenever one of the
arguments $A_i$

